\section{Project Structure}

%Definisci gli Actors: Replica, Client e altri che possono essere utili dopo
%Tabella/lista dei message types, potremmo separali per fase (oppure no, dipende dallo spazio che abbiamo)

\subsection{Actors}

The system is implemented as a set of Akka actors, all deriving from the abstract base classes
\texttt{AbstractReplica} and \texttt{AbstractClient} provided by the codebase, which expose the
mandatory API callbacks and the crash/initialization message handling shared by every node. All
protocol state (position list, update log, epoch info, pending timeouts) lives inside a single
\texttt{Replica} actor instance and is never exposed by reference outside it; every message crossing an actor boundary
carries only immutable data. As a result, no thread-level synchronization primitive is ever required within a single replica, a property of the actor
model, not specific to our design. 
%This is unrelated to \emph{distributed} synchronization between replicas, which is instead achieved explicitly through the update and election protocols.

\begin{itemize}[leftmargin=*]
  \item \textbf{Replica}: implements a full storage node, running the three protocols (update broadcast, heartbeat-based failure detection, ring election) as
  message handlers over its own private state. Concentrating all three protocols in a single actor
  class avoids introducing extra inter-actor
  coordination for state that all three protocols must read consistently, such as the current epoch
  and coordinator id.

  \item \textbf{Client}: issues \texttt{ReadRequest}/\texttt{WriteRequest} messages to a replica and
  tracks each outstanding request together with its own timeout, invoking the corresponding callback
  when a result, or a timeout, is received. Requests are tracked individually so a slow read never blocks an unrelated write to a different index.

  \item \textbf{NetworkChannel}: a per-destination actor, transparently instantiated by the
  \texttt{tell()} helper in \texttt{AbstractReplica}, that injects a random delay between
  \texttt{minLatency} and \texttt{maxLatency} before forwarding a message to its real recipient. One
  channel is created per (sender, destination) pair specifically so FIFO ordering is preserved between
  any two actors.

\end{itemize}

\subsection{How Updates Are Stored}

Every accepted update is uniquely identified by an \texttt{UpdateTimestamp}, the $\langle e,i
\rangle$ pair (epoch, sequence number) that also defines the total order between updates via
\texttt{compareTo}. Three data structures build on this identifier to track an update through its
lifecycle:

\begin{itemize}[leftmargin=*]
  \item \textbf{PendingUpdate}: created by the coordinator when a broadcast round starts. It bundles
  the update data, its timestamp, the client to reply to, and the set of replica ids that have ACKed
  it so far, so quorum can be checked with a simple size comparison instead of re-scanning messages.
  \item \textbf{UpdateLog}: the durable, per-replica history of \emph{completed} updates. Entries are
  appended in strict timestamp order (\texttt{addLog}) once a \texttt{WriteOkMsg} is processed; this
  is also what lets a newly elected coordinator determine "the most recent update it knows of" during
  the election. During synchronization, out-of-order insertion (\texttt{addLogIfAbsent}) is used
  instead, since a recovered update may need to be inserted before entries a replica already applied.
  \item \textbf{PositionList} / \textbf{PersonOfInterest}: the actual application state (the array of
  positions), where each entry additionally remembers the timestamp of the update that last wrote it.
  This is only ever touched by \texttt{applyUpdate()}, after an update has already been safely logged.
\end{itemize}

Keeping "in-progress" state (\texttt{PendingUpdate}) and "completed" state (\texttt{UpdateLog}) in
two distinct structures makes it straightforward, during recovery, to tell whether an interrupted
update was merely broadcast or already confirmed by a quorum.

\subsection{Message Flow}

\begin{center}
\small
\begin{tabular}{|C{2.8cm}|C{9.3cm}|}
\hline
\textbf{Phase} & \textbf{Purpose} \\
\hline
Client $\leftrightarrow$ Replica &
  Clients issue reads (answered immediately from local state) and writes (forwarded to the
  coordinator if needed); timeouts on the client side detect an unresponsive replica. \\
\hline
Update broadcast &
  Implements the two-phase \texttt{UPDATE}/\texttt{ACK}/\texttt{WRITEOK} exchange coordinated by the
  coordinator, plus the timeouts used to detect a coordinator crash mid-round. \\
\hline
Heartbeat &
  Periodic coordinator-to-replica liveness check, used both to let replicas detect a silently
  crashed coordinator and to let the coordinator remove unresponsive replicas. \\
\hline
Coordinator election &
  Ring-based election of a new coordinator after a crash is detected, followed by a synchronization
  phase. \\
\hline
Setup / test instrumentation &
  System initialization (\texttt{InitSystem}), crash injection (\texttt{Crash}), and the
  observability callbacks (\texttt{CoordinatorElected}, \texttt{UpdateApplied},
  \texttt{ElectionStarted}) for the automated test suite.\\
\hline
\end{tabular}
\end{center}